\par\Needspace{6\baselineskip}\medskip
\hypertarget{native-081-article062}{}\pdfbookmark[2]{62.}{article-062}
\begin{center}62.\end{center}
Let the celestial body have heliocentric longitude $\lambda$ and latitude $\beta$, geocentric longitude $l$ and latitude $b$, and distances $r$ from the Sun and $\Delta$ from the Earth. Finally, let the Earth have heliocentric longitude $L$ and latitude $B$, and distance $R$ from the Sun. Since we do not assume $B=0$, our formulas also apply when the heliocentric positions and the geocentric position are referred to any other plane in place of the ecliptic. We need only drop the names longitude and latitude. Moreover, parallax can be taken into account at once, provided only that the heliocentric position of the Earth is referred directly to a point on its surface in place of its center. Let us further set $r\cos\beta=r'$, $\Delta\cos b=\Delta'$, $R\cos B=R'$. Now refer the positions in space of the celestial body and the Earth to three planes, of which one is the ecliptic and the second and third have their poles at longitudes $N$ and $N+90^0$. The following equations result at once
\[
\begin{gathered}
r'\cos(\lambda-N)-R'\cos(L-N)=\Delta'\cos(l-N)\\
r'\sin(\lambda-N)-R'\sin(L-N)=\Delta'\sin(l-N)\\
r'\operatorname{tang}\beta-R'\operatorname{tang}B=\Delta'\operatorname{tang}b,
\end{gathered}
\]
where the angle $N$ is entirely arbitrary. The first and second equations immediately give $l-N$ and $\Delta'$, from which the third yields $b$. From $b$ and $\Delta'$ we obtain $\Delta$. To make the calculation as convenient as possible, we now choose the arbitrary angle $N$ in the following three ways.

I. Setting $N=L$, we put $\frac{r'}{R'}\sin(\lambda-L)=P$ and $\frac{r'}{R'}\cos(\lambda-L)-1=Q$. We then find $l-L$, $\frac{\Delta'}{R'}$, and $b$ by the formulas
\[
\begin{gathered}
\operatorname{tang}(l-L)=\frac{P}{Q}\\[.5em]
\frac{\Delta'}{R'}=\frac{P}{\sin(l-L)}=\frac{Q}{\cos(l-L)}\\[.5em]
\operatorname{tang}b=\frac{\frac{r'}{R'}\operatorname{tang}\beta-\operatorname{tang}B}{\frac{\Delta'}{R'}}\mathord{\cdot}
\end{gathered}
\]

\SourceLeaf{082}II. Setting $N=\lambda$, we put $\frac{R'}{r'}\sin(\lambda-L)=P$ and $1-\frac{R'}{r'}\cos(\lambda-L)=Q$. We then have
\[
\begin{gathered}
\operatorname{tang}(l-\lambda)=\frac{P}{Q}\\[.5em]
\frac{\Delta'}{r'}=\frac{P}{\sin(l-\lambda)}=\frac{Q}{\cos(l-\lambda)}\\[.5em]
\operatorname{tang}b=\frac{\operatorname{tang}\beta-\frac{R'}{r'}\operatorname{tang}B}{\frac{\Delta'}{r'}}\mathord{\cdot}
\end{gathered}
\]

III. Setting $N=\frac12(\lambda+L)$, we find $l$ and $\Delta'$ by the equations
\[
\begin{gathered}
\operatorname{tang}\{l-\tfrac12(\lambda+L)\}=\frac{r'+R'}{r'-R'}\operatorname{tang}\tfrac12(\lambda-L)\\[.5em]
\Delta'=\frac{(r'+R')\sin\frac12(\lambda-L)}{\sin(l-\frac12(\lambda+L))}=\frac{(r'-R')\cos\frac12(\lambda-L)}{\cos(l-\frac12(\lambda+L))}
\end{gathered}
\]
and then $b$ by the equation given above. The logarithm of the fraction $\frac{r'+R'}{r'-R'}$ is readily calculated by setting $\frac{R'}{r'}=\operatorname{tang}\zeta$, which gives $\frac{r'+R'}{r'-R'}=\operatorname{tang}(45^0+\zeta)$. In this form method III is somewhat shorter than I and II for determining $l$. For the remaining operations, however, we consider I and II preferable to III.
